% 3. วิธีดำเนินงาน -- ย่อจากรายงานบทที่ 3 (../report/chapter/03-methodology.tex)
\section{วิธีดำเนินงาน}

\subsection{กลุ่มตัวอย่างและเครื่องมือ}

ผู้ตอบคือนักศึกษาระดับปริญญาตรี มหาวิทยาลัยขอนแก่น เลือกแบบตามสะดวก (convenience sampling)
โดยส่งลิงก์แบบสอบถามในกลุ่มออนไลน์ของคณะและชั้นปีต่าง~ๆ
ระหว่างวันที่ 23 สิงหาคม ถึง 3 กันยายน 2569 ได้คำตอบ 121 ชุด
แบบสอบถามเป็น Google Forms 36 ข้อ ประกอบด้วยข้อมูลทั่วไป พฤติกรรมการใช้โซเชียลมีเดีย
การนอน การออกกำลังกาย ความกดดันจากการเรียน เกรดเฉลี่ยสะสม ความเพียงพอทางการเงิน
ข้อดักความใส่ใจ 1 ข้อ และแบบวัด DASS-21 ฉบับภาษาไทย \cite{webster2008thai} โดยไม่แก้ข้อความ
ข้อชั่วโมงการใช้งานให้ผู้ตอบเปิดดูค่าที่มือถือบันทึกไว้ (Screen Time หรือ Digital Wellbeing)
แทนการประมาณจากความจำ

ข้อดักความใส่ใจ (attention check) ขึ้นต้นว่า ``ข้อนี้ไม่ใช่คำถาม''
และสั่งให้เลือกตัวเลือกที่ 3 จาก 4 ตัวเลือก ผู้ที่เลือกตัวเลือกอื่นถือว่าไม่ได้อ่านคำถามจริง
ผู้ตอบทุกคนต้องยืนยันความยินยอมก่อนตอบ แบบสอบถามไม่เก็บข้อมูลที่ระบุตัวผู้ตอบได้
และแสดงสายด่วนสุขภาพจิต 1323 ทั้งตอนต้นและตอนท้าย

\subsection{ตัวแปร}

ตาราง~\ref{tab:variables} แสดงตัวแปรเชิงปริมาณที่ใช้ในการวิเคราะห์ พร้อมค่าเฉลี่ยและส่วนเบี่ยงเบนมาตรฐาน
นอกจากนี้ยังมีประเภทเนื้อหาที่ดูมากที่สุด ซึ่งมี 8 ตัวเลือกในแบบสอบถาม
และยุบเหลือ 4 กลุ่ม คือ ข่าว บันเทิง ชีวิตคนอื่น และอื่น~ๆ ตามกติกาที่กำหนดไว้ก่อนเห็นผล
คะแนน DASS-21 เป็นแบบวัดเชิงลบ คะแนนยิ่งสูงยิ่งมีอาการมาก
ส่วนคุณภาพการนอนและความเพียงพอทางการเงิน คะแนนยิ่งสูงยิ่งดี

\begin{table}[!t]
\centering\small
\caption{ตัวแปรที่ใช้ในการวิเคราะห์และสถิติเชิงพรรณนา ($n = 107$)}
\label{tab:variables}
\begin{tabularx}{\columnwidth}{|L|c|r|r|}
\hline
\textbf{ตัวแปร} & \textbf{ช่วงค่า} & \textbf{เฉลี่ย} & \textbf{SD} \\
\hline
\multicolumn{4}{|l|}{\textit{ตัวแปรตาม}} \\
\hline
ภาวะซึมเศร้า & 0--42 & 9.42 & 8.13 \\
ความวิตกกังวล & 0--42 & 11.40 & 7.47 \\
ความเครียด & 0--42 & 11.89 & 7.78 \\
\hline
\multicolumn{4}{|l|}{\textit{ตัวแปรต้นหลัก}} \\
\hline
ชั่วโมงใช้โซเชียลมีเดียต่อวัน & 0--24 & 7.83 & 2.70 \\
วันที่ติดตามข่าวเชิงลบใน 7 วัน & 0--7 & 3.46 & 1.54 \\
\hline
\multicolumn{4}{|l|}{\textit{ตัวแปรควบคุม}} \\
\hline
ชั่วโมงการนอนต่อคืน & 0--24 & 6.21 & 1.54 \\
คุณภาพการนอน & 1--10 & 5.54 & 2.16 \\
วันที่ออกกำลังกายต่อสัปดาห์ & 0--7 & 2.00 & 1.81 \\
ความกดดันจากการเรียน & 1--10 & 6.12 & 2.26 \\
เกรดเฉลี่ยสะสม & 0--4 & 3.20 & 0.49 \\
ความเพียงพอทางการเงิน & 1--10 & 6.50 & 2.13 \\
\hline
\end{tabularx}

\vspace{2pt}
\raggedright หมายเหตุ: ชั่วโมงการนอน $n = 106$ และเกรดเฉลี่ยสะสม $n = 104$ เนื่องจากมีค่าว่าง
\end{table}

\subsection{การทำความสะอาดข้อมูล}

คำตอบที่เป็นข้อความปนตัวเลข เช่น ``6-7 ชม.'' แปลงเป็นค่าเฉลี่ยของตัวเลขในช่อง
ค่าที่เป็นไปไม่ได้ เช่น นอน 210 ชั่วโมง หรือเกรดเฉลี่ยเกิน 4.00 กำหนดเป็นค่าว่าง
ผู้ที่ไม่ผ่านข้อดักความใส่ใจ 14 คน (ร้อยละ 11.6) ถูกตัดออกทั้งแถว เหลือ 107 คน
ค่าว่างในข้อที่ไม่บังคับตอบไม่ได้เติมค่า แต่ละการวิเคราะห์ใช้แถวที่มีข้อมูลครบเฉพาะตัวแปรที่ใช้
การถดถอยพหุคูณจึงเหลือ 103 คน
คะแนนและระดับความรุนแรงของ DASS-21 คำนวณตามเกณฑ์ในคู่มือ \cite{lovibond1995manual}

\subsection{การวิเคราะห์ข้อมูล}

ความสัมพันธ์ระหว่างพฤติกรรมการใช้โซเชียลมีเดียกับคะแนนแต่ละมิติวัดด้วยสหสัมพันธ์เพียร์สัน (\stat{r})
พร้อมช่วงความเชื่อมั่น 95\% แบบ bootstrap 5,000 รอบ
และตีความขนาดตามเกณฑ์ของ Cohen \cite{cohen1988statistical}
คือ 0.10 เล็ก 0.30 ปานกลาง และ 0.50 ใหญ่
สหสัมพันธ์บางส่วน (partial correlation) ใช้ดูความสัมพันธ์หลังหักผลของชั่วโมงการนอน
คุณภาพการนอน ความกดดันจากการเรียน เกรดเฉลี่ย และความเพียงพอทางการเงิน
การเทียบคะแนนระหว่างกลุ่มใช้การทดสอบ Kruskal-Wallis และการวิเคราะห์ความแปรปรวนทางเดียว (ANOVA)

คำถามว่าความสัมพันธ์ยังเหลืออยู่หรือไม่เมื่อควบคุมปัจจัยอื่น
ตอบด้วยการถดถอยพหุคูณแบบลำดับขั้น (hierarchical regression)
ขั้นที่ 1 ใส่ตัวแปรควบคุมหกตัว ขั้นที่ 2 เพิ่มชั่วโมงใช้งานและวันที่ติดตามข่าว
แล้วดูค่า \Rsq{} ที่เพิ่มขึ้น ($\Delta$\Rsq) ด้วยการทดสอบ F-change
สมการเต็มของแต่ละมิติเป็นดังสมการ~(\ref{eq:model})
%
\begin{equation}
    Y_k = \beta_0 + \beta_1 X_1 + \beta_2 X_2 + \cdots + \beta_8 X_8 + \varepsilon
    \label{eq:model}
\end{equation}
%
เมื่อ $Y_k$ คือคะแนนมิติที่ $k$ $X_1$ คือชั่วโมงใช้งานต่อวัน $X_2$ คือวันที่ติดตามข่าวเชิงลบ
และ $X_3$ ถึง $X_8$ คือตัวแปรควบคุมหกตัวในตาราง~\ref{tab:variables}
ค่าความคลาดเคลื่อนมาตรฐานของสัมประสิทธิ์ใช้แบบ HC3
เนื่องจากคะแนน DASS-21 เบ้ขวาและความแปรปรวนของค่าคลาดเคลื่อนอาจไม่คงที่
ค่า \stat{p} ของการทดสอบชุดหลักปรับด้วยวิธี Bonferroni (คูณ 3 ตามจำนวนสมการ)
และตรวจข้อตกลงเบื้องต้นด้วย VIF, residual plot, Q-Q plot, การทดสอบ Breusch-Pagan
และ Shapiro-Wilk
การวิเคราะห์ทั้งหมดใช้ Python ร่วมกับไลบรารี pandas, scipy และ statsmodels
ที่ระดับนัยสำคัญ 0.05
